\chapter{Conclusion}\label{ch:conclusion}

In the preceding chapter we evaluated each optimisation level in isolation and then measured their combined effect, revealing that tile shaving and columnar re-encoding dominate the load-time improvements while style-level passes provide modest, additive contributions.
In this chapter we draw together the findings: we summarise our contributions, distil the key insights, revisit the three research questions posed in \cref{sec:research_goals}, and identify directions for future work.

\section{Summary of Contributions}

Our main contributions are:

\begin{enumerate}
    \item We present a \textbf{three-level optimisation pipeline} (\cref{sub:pipeline}) that proved most effective as an enabler: expression-level simplifications cascaded into forms that structural passes could exploit, with constant folding and dead elimination together accounting for the largest per-pass rendering gains (\cref{sub:combined_marginal}).
    We observed that some passes, such as the default removal passes, had both positive and negative effects on the metrics measured, but not to a significant enough extent to alter the overall trend.

    \item We show that \textbf{style-driven tile shaving} (\cref{sub:tile_shaving}) is the dominant contributor to load-time and transfer-size improvements.  Dead elimination narrows the pruning advisory, enabling elimination of property columns and tiles from the encoded output, though we find the interaction with style-level passes to be additive rather than synergistic (\cref{sub:mlt_discussion}).

    \item We contribute \textbf{automatic encoding strategy selection} for the \ac{MLT} format, which closed the gap that the fixed-plan Java encoder left open: the Rust encoder's per-stream profiling and sort-strategy competition reduced tile volume by \qty{28}{\percent} over the Java baseline, and we found this advantage to persist under every wire compressor tested (\cref{tab:mvt_mlt_comparison}).

    \item Through our \textbf{empirical evaluation} (\cref{ch:experiments}) we confirm that the three optimisation categories are complementary: expression and structural passes dominate rendering-throughput improvements, while data-level passes dominate load-time improvements, with the full pipeline reducing total tile volume by \qty{28}{\percent} (\qty{12}{\percent} over gzip-compressed \ac{MVT}) and improving end-to-end load times by \qtyrange{71}{75}{\percent} on the two styles we benchmarked with tile data.
\end{enumerate}

\section{Key Findings}

From our evaluation we draw three architectural insights that go beyond the quantitative results we report in \cref{ch:experiments}.

\emph{First}, we find that tile shaving dominates the load-time and transfer-size improvements (\cref{sub:mlt_discussion}).
Style-level passes simplify the advisory computation and provide modest rendering improvements, but our aggregate data shows no super-additive interaction: the combined style+shaving reduction is approximately the sum of the individual reductions.
For the geospatial community this means that tile shaving and encoding strategy selection are the primary levers; style-level passes are useful enablers that improve code quality and provide small rendering gains, rather than co-equal partners in a synergistic pipeline.
For FPS (\qtyrange{162}{213}{\percent}) and transfer size (\qty{36}{\percent}) the full pipeline's effect is significant, while the incremental contribution of style-level passes over shaving alone is modest.

\emph{Second}, we identify three plausible improvement channels (inferred from proxy metrics, not direct per-mechanism instrumentation) with deployment implications: structural passes dominate rendering-throughput improvements, data-level passes dominate load-time improvements, and expression passes serve primarily as enablers.
A mobile deployment on a constrained network should prioritise tile shaving and re-encoding, whereas a desktop deployment bottlenecked by draw calls benefits most from structural passes alone.

\emph{Third}, we show that the compiler-optimisation framing - peephole rewriting, fixpoint iteration, dead code elimination - transfers effectively to the map-style domain.
Our rewrite rules (\cref{tab:all_rewrite_rules}) are individually simple, yet their composition through fixpoint iteration achieves globally significant simplification across styles of widely varying complexity (raw reductions of \qtyrange{10}{71}{\percent} across the Maputnik generalisation corpus).
We attribute this transfer to a key structural property shared by both domains: expressions are compositional (sub-expressions can be rewritten independently without invalidating the enclosing context), and the rewrite rules are semantics-preserving, so fixpoint iteration over a monotone rule set converges to a sound result regardless of rule application order.

\section{Research Questions Revisited}\label{sec:rq_revisited}

We can now answer the three research questions posed in \cref{sec:research_goals} directly from the experiments in \cref{ch:experiments}.

\paragraph{R1 - To what extent can automatic data and style co-optimisations reduce map tile size and improve rendering performance?}
We resolve the $(S, T) \to (S', T')$ optimisation of \cref{sec:problem_statement} empirically.
Our style-level evaluation across all 15 styles and 18 scenarios shows style-gzip reductions of \qtyrange{3}{47}{\percent} (median ${\sim}$\qty{8}{\percent}).
At the tile-data level, our full-pipeline benchmarks on the two styles we evaluated end-to-end (Fiord and Liberty) show a total tile volume reduction of \qty{28}{\percent} (from \qty{4.24}{\giga\byte} to \qty{3.06}{\giga\byte} on the Germany OpenMapTiles tileset), \qty{12}{\percent} over gzip-compressed \ac{MVT}, and load-time improvements of \qtyrange{71}{75}{\percent} (\cref{sub:combined_ablation}).
We find the gains are not uniform: verbose institutional styles with hundreds of layers benefit far more than already-compact basemaps, and scenarios dominated by overview zoom levels benefit more than those dominated by street-level panning (\cref{sub:combined_marginal}).
We therefore answer R1 with \enquote{yes, meaningfully, but in a style- and scenario-dependent way}: tile shaving and the Rust encoder's strategy competition account for most of the tile-size and load-time reduction, while dead elimination plus layer merging account for most of the frame-rate improvements.  We establish the load-time claim of \qtyrange{71}{75}{\percent} for only two styles (Fiord and Liberty); we evaluate the remaining styles at the style level only, and tile-shaving effectiveness varies from \qty{19.8}{\percent} to \qty{46.3}{\percent} across styles.

\paragraph{R2 - Which techniques are most effective, and how do they interact?}
Our cumulative ablation in \cref{sub:combined_ablation} and per-style results in \cref{sec:app_per_style} rank the passes by marginal contribution and reveal three clusters.
\emph{First}, we find that constant folding (including stats-driven folding) and expression simplification dominate the expression-level reductions because they simplify filters into forms that the structural passes can exploit, and because a single successful fold cascades through enclosing expressions.
\emph{Second}, we observe that dead elimination and layer merging dominate the structural-level reductions because they reduce the per-frame draw-call count directly; we also find them to be strongly synergistic, because dead elimination removes gaps between mergeable layers that would otherwise prevent the merge pass from recognising adjacency.
\emph{Third}, data-level optimisations operate on a different axis: tile shaving reduces transfer size by using the optimised style's property and value references to build a pruning advisory, and the \ac{MLT} encoder's sort-strategy and per-stream competition reduces encoded size on top of that.
Style-level dead elimination narrows the pruning advisory, which can shrink tiles further, but our aggregate data shows this interaction is additive rather than synergistic.

\paragraph{R3 - How expensive are these optimisations in terms of preprocessing time, and what are the tradeoffs between decoding latency, rendering throughput, and transfer size?}
Our preprocessing-cost model in \cref{sub:preprocessing_cost} decomposes the total cost into four additive stages (style optimisations, statistics collection, advisory computation, \ac{MLT} re-encoding) and identifies re-encoding as the dominant term.
Our empirical measurement on the Germany OpenMapTiles tileset (\cref{tab:preprocessing_cost}) shows that the multi-strategy competition increases single-threaded re-encoding cost by a factor of $1.95\times$ over a fixed-plan encoder - well below the na\"ive $4\times$ expectation, thanks to bounding-box pruning, data profiling, and warm-start caching.
The full single-threaded pipeline completes in under \qty{7}{\minute}; with 16-thread parallelism the wall-clock time drops to approximately \qty{2}{\minute}, paid once per tileset build and amortised over every subsequent tile request.
On the tradeoff axis, we find that decoding latency improvement is the most dramatic gain: the \ac{MLT}-Rust decoder achieves \numrange{4.6}{8.3}$\times$ higher full-decode throughput than \ac{MVT}+gzip across zoom levels 4, 7, and 13 (\cref{tab:decode_throughput}).
Transfer-size reduction is the most reliably measurable axis (deterministic under compression), but for style-level passes it primarily affects initial load; tile-level volume reduction (\qty{12.4}{\percent} for \ac{MLT}-Rust vs.\ \ac{MVT} under gzip) has the larger sustained impact as tiles are fetched continuously.
Rendering throughput (FPS) for style-only optimisations varies by style (median change ${\sim}$\qty{3}{\percent}, range \qtyrange{-4}{13}{\percent}), affected almost entirely by the structural passes rather than by the data-level passes - confirming that the multi-level pipeline is a set of complementary mechanisms that each address a different bottleneck, with tile shaving carrying the dominant share of load-time improvement.
We do not measure energy consumption directly (our benchmark host is a desktop workstation rather than a battery-constrained device), but the transfer-size and frame-rate results are well-established proxies on mobile targets and we translate them into the expected battery-life direction in the qualitative discussion in \cref{sub:combined_discussion}.

\section{Future Work}\label{sec:conclusion_future_work}

We see several directions that remain open:

\begin{itemize}
    \item \textbf{Symbol layer merging}: we exclude symbol layers from merging in our current implementation because it would disrupt collision detection. We leave investigating collision-aware merging strategies to future work, as it could yield additional layer count reductions on label-heavy styles.
    \item \textbf{Runtime re-encoding}: our current tile encoding pipeline is offline. Exploring on-the-fly MVT-to-MLT re-encoding at the tile server or \ac{CDN} edge could make the encoding optimisations available without a full tile rebuild. The AnyBlox framework~\cite{gienieczkoAnyBloxFrameworkSelfDecoding2025}, which bundles lightweight WebAssembly decoders with datasets, suggests an alternative path: bundling an \ac{MLT} decoder as a WASM module alongside the tile data would decouple format evolution from client updates.
    \item \textbf{Broader style coverage}: we leave extending the evaluation to proprietary and commercial styles (Mapbox Streets, Google Maps style equivalents) to future work, as it would strengthen the generalisability claims.
    \item \textbf{Client-side integration}: we plan to integrate the optimised \ac{MLT} decoder into MapLibre GL~JS and measure end-to-end improvements in a production deployment.
    \item \textbf{Selective column decoding}: the \ac{MLT} columnar layout permits skipping unreferenced columns entirely during decoding - a form of \emph{late materialisation}~\cite{abadiMaterializationStrategiesColumnOriented2007} applied at the tile level. If the decoder were augmented with a column projection derived from the optimised style's property references, decode cost would scale with the number of columns actually used rather than with the total schema width - an optimisation complementary to tile shaving, which removes columns at encoding time (early materialisation) rather than skipping them at decode time.
    \item \textbf{Fractional-zoom visual correctness testing}: our current pixelmatch harness tests integer zoom levels only. Metadata refinement and zoom-bounded ramp folding operate on fractional zoom thresholds, and bugs are most likely to appear between integer zooms. We leave adding fractional-zoom samples (e.g., $z + 0.25$, $z + 0.5$, $z + 0.75$) to future work, as they would more completely sample the $\forall v \in \mathcal{V}$ quantifier of \cref{eq:visual_equivalence} and thus strengthen the soundness argument of \cref{sec:problem_statement}.
    \item \textbf{MinHash threshold sensitivity}: we tuned the MinHash similarity threshold ($\tau = 7.5\%$) on the Germany OpenMapTiles tileset (\cref{fig:minhash_sweep}), but we have not characterised sensitivity across tilesets with different vocabulary distributions. Repeating the sweep on non-OSM schemas would validate the choice.
    \item \textbf{Serving-time applicability}: our current pipeline assumes offline batch processing, which is impractical for dynamic tile sources (e.g., PostGIS-backed tile servers with frequently updated data).
    Not all passes require offline processing, however.
    The style optimiser operates on the style alone and can run at deployment time (when the style is published) rather than at tile-build time.
    The tile pruning advisory could potentially be applied as a lightweight per-request filter at serving time, since it reduces to column selection and predicate evaluation - operations that tile servers already perform for other purposes - connecting directly to the \emph{adaptive query processing} literature~\cite{deshpandeAdaptiveQueryProcessing2006}, which studies techniques for adjusting query plans at runtime in response to changing data characteristics.
    The encoding strategy competition, by contrast, requires access to the full column to profile and is inherently a batch operation; its results could be cached per schema and reused across tile requests.
    We leave characterising which passes are applicable at serving time and measuring their overhead on a request-critical path to future work, as it would extend the pipeline's applicability to dynamic tile sources.
    \item \textbf{Learned sort-strategy prediction}: our encoder's sort-strategy competition (\cref{sub:data_optimisations}) exhaustively trials Hilbert, Morton, ID, and unsorted orderings and retains the smallest encoding. Kraska et~al.~\cite{kraskaCaseLearnedIndex2017} showed that learned models can replace traditional index structures by predicting key positions; by analogy, a learned model trained on layer-level statistics (cardinality, spatial spread, value distribution) could predict the optimal sort strategy without running all trials, reducing preprocessing cost while preserving most of the size reduction.
    \item \textbf{Equality saturation for the rewrite engine}: our current rewrite system applies rules in a fixed priority order and is not confluent (\cref{sub:fixpoint}). Replacing the sequential rewriter with an equality saturation engine~\cite{tateEqualitySaturationNew2010,willseyEggFastExtensible2020} would explore all equivalent forms simultaneously and select the globally optimal one, eliminating the ordering dependence. Beyond execution, Nandi et~al.~\cite{nandiRewriteRuleInference2021} showed that equality saturation can also be used to \emph{infer} rewrite rules automatically from input--output examples, which could complement or replace the manually authored rules we present in this thesis. We see this direction as primarily valuable if the rule set grows significantly or if ordering sensitivity becomes a practical concern.
\end{itemize}
